\section{Introduction}
\label{sec:introduction}

A smooth cubic surface has \(27\) geometric lines, but the base field need
not see any of them individually.  Their Galois action can split into several
orbits, and even a transitive action need not be the full incidence group
\(W(E_6)\).  For a fixed surface, determining the line field therefore
requires more than smoothness or the classical line count.

This paper carries out that determination for the exact cubic surface
\(Y/\Q\) displayed in \cref{eq:cubic}.  The equation arose as a
four-variable Yukawa cubic in the preceding H\'enon-core construction, but
the present problem is independent of the Hodge-theoretic origin.  Its object
is the finite scheme
\[
  \FanoY\subset\Gr(2,4),
\]
and its central questions are arithmetic: is this scheme connected, what is
the normal field of its geometric points, and what does the resulting action
fix in the Picard lattice?

Three points make the calculation delicate.  First, a zero-dimensional
calculation on one Grassmann chart does not automatically recover the global
Fano scheme.  We close that gap by first proving that \(\FanoY\) is finite
\(\etale\) of rank \(27\); an open chart is then open-and-closed, and a
rank-\(27\) closed subscheme must be the whole line scheme.  Second, a
transitive subgroup containing an element of order five may still be the
index-two simple subgroup \(U\subset W(E_6)\).  The required parity is the
determinant in the \(E_6\) reflection representation, not ordinary
permutation sign: every element of \(W(E_6)\) acts evenly on the \(27\)
lines.  Third, the degree-\(27\) residue field of one line is not the common
normal line field.  We denote the former by \(E\) and its normal closure by
\(K\) throughout.

The main result is as follows.

\begin{theorem}[Line field, Galois closure, and Picard ranks]
\label{thm:main}
For the cubic surface \(Y/\Q\) in \cref{eq:cubic}, there is an irreducible
polynomial \(g\in\mathbf Z[d]\) of degree \(27\) such that
\[
  \FanoY\cong\Spec(E),\qquad E=\Q[d]/(g).
\]
Thus \(\FanoY\) is connected and finite \(\etale\) of rank \(27\).
If \(K\) is the splitting field of \(g\), then
\[
  \Gal(K/\Q)\cong W(E_6),\qquad [K:\Q]=51840.
\]
Moreover,
\[
  \rho(Y_{\Qbar})=7,\qquad \rho(Y/\Q)=1.
\]
If a finite extension \(L/\Q\) defines a line on \(Y\), then
\[
  27\mid[L:\Q].
\]
In particular, \(Y\) has no \(\Q\)-rational line.
\end{theorem}

The proof combines written geometry and group theory with finite exact
calculations.  Its four principal contributions are:
\begin{enumerate}[leftmargin=*]
\item We reconstruct the complete line scheme from one exact Grassmann chart.
The degree-order quotient has dimension \(27\) and Hilbert counts
\((1,4,10,12,0)\); a four-polynomial lexicographic shape gives a
degree-\(27\) eliminant and three linear back-substitutions.  All four
original line equations reduce to zero.
\item We prove irreducibility by complete modular data rather than a
black-box rational factorization.  The factor-degree subset-sum intersections
shrink through
\[
\{0,3,6,\ldots,27\},\quad
\{0,6,9,12,15,18,21,27\},\quad
\{0,9,18,27\},\quad
\{0,27\}.
\]
\item We prove maximal Galois action.  The prime \(37\) gives cycle type
\((2,5,5,5,10)\).  Exact enumeration finds \(5184\) elements of this type in
\(W(E_6)\), all outside the Coxeter-even subgroup \(U\), and none inside it.
Together with the transitive order-five criterion of
Elsenhans--Jahnel, this forces the full Weyl group
\cite{ElsenhansJahnel2009}.
\item We reconstruct the Picard reflection model and find a
one-dimensional fixed space.  Hochschild--Serre is used only after tensoring
with \(\Q\), so the conclusion is equality of ranks rather than an
unqualified integral equality of Picard groups.
\end{enumerate}

The arithmetic conclusion about lines is intentionally narrow.  A surface
without a rational line may still have rational points, and nothing here
decides rationality, stable rationality, the Hasse principle, or a
Brauer--Manin obstruction.  The Yukawa origin supplies the distinguished
equation but does not turn the line-field computation into a motivic or
Calabi--Yau realization theorem.

The rest of the paper is organized as follows.
\Cref{sec:surface-fano} fixes the surface and establishes the global finite
\(\etale\) line scheme.  \Cref{sec:chart-scheme} gives the chart presentation
and the open-and-closed bridge.  \Cref{sec:irreducibility} proves that the
degree-\(27\) algebra is a field.  \Cref{sec:weyl-galois} determines its
normal closure, and \cref{sec:picard-lines} derives the Picard and
line-definition consequences.  Exact replay and scope are recorded in
\cref{sec:exact-replay,sec:context-scope}; the appendices contain source
locators and the compact exact data.
